% GENERATED FILE. Edit canonical lessons, then run npm run generate:books.
% canonical-source-hash: fnv1a64:9ebe2c7b564d9325
% canonical-lessons: LA-C449-vappa, LA-C449-merum, LA-C449-diluo, LA-C449-putidus, LA-C449-copa

\chapter{Flat wine, Neat wine, To water down, Rotten, Tavern hostess}
\label{ch:la-flat-wine-neat-wine-to-water-down-rotten-tavern-hostess}
\hlchaptermodality{449}

\begin{chapteropening}
\textbf{By the end of this chapter:} \emph{I can complain about the food and drink at an inn: flat wine, neat wine, watered wine, rotten food, and the tavern hostess.}

Complete the last lesson of chapter 449: I can complain about the food and drink at an inn: flat wine, neat wine, watered wine, rotten food, and the tavern hostess.
\end{chapteropening}

\section[vappa, vappae]{vappa, vappae — flat wine, wine gone off; a good-for-nothing}

\label{lesson:LA-C449-vappa}



\begin{quote}
Before the new one: say the Latin for a drinker, then the Latin for hoarse.
\end{quote}

\subsection*{You'll want to know: vappa, vappae}
\textbf{vappa, vappae} — \textquotedblleft{}flat wine, wine gone off; a good-for-nothing\textquotedblright{}.

\begin{quote}
Caupō nōn vīnum bonum sed vappam vendit.
\end{quote}

\begin{quote}
\emph{The innkeeper sells not good wine but flat wine.}
\end{quote}

First declension, feminine: \emph{vappa, vappae f.} Wine that has lost its taste; also, of a person, a worthless fellow.

\begin{cousinweb}
\emph{vappa} is usually linked with \emph{vapor} 'steam' (Lewis \& Short): wine whose spirit has gone up in vapour. Lewis \& Short also compare \emph{vapidus} 'flat', the source of English \emph{vapid}.
\end{cousinweb}

\begin{grammarlens}[title={What you've built}]
One more word for telling, explaining and describing at B1.
\end{grammarlens}

\subsection*{Your turn}
\begin{itemize}
\raggedright
  \item \emph{Say it:} \emph{vappa, vappae}
  \item \emph{Say it:} \emph{vappa, vappae}, once more
  \item \emph{Say it:} a sentence of your own with \emph{vappa, vappae}
  \item \emph{Recall:} say the Latin for a drinker, then the Latin for hoarse, then say \emph{vappa, vappae} again
\end{itemize}

\begin{tcolorbox}[breakable,colback=teal!4,colframe=teal!35!black,title={Before you move on}]
What does vappa, vappae mean? (\textquotedblleft{}Flat wine, wine gone off; a good-for-nothing\textquotedblright{}.) Say it once more.
\end{tcolorbox}
\section[merum, merī]{merum, merī — neat wine, wine without water}

\label{lesson:LA-C449-merum}



\begin{quote}
Before the new one: say the Latin for hoarse, then the Latin for flat wine.
\end{quote}

\subsection*{You'll want to know: merum, merī}
\textbf{merum, merī} — \textquotedblleft{}neat wine, wine without water\textquotedblright{}.

\begin{quote}
Pōtor merum bibit, nōn vīnum cum aquā.
\end{quote}

\begin{quote}
\emph{The drinker drinks neat wine, not wine with water.}
\end{quote}

Second declension, neuter: \emph{merum, merī n.}, the neuter of \emph{merus} 'pure, unmixed'. Romans mixed wine with water; drinking it neat was a sign of excess. English \emph{mere} comes from \emph{merus}.

\begin{cousinweb}
\emph{merum} is the neuter of the adjective \emph{merus} 'pure, unmixed' (Lewis \& Short), used as a noun for wine with no water. English \emph{mere} comes from \emph{merus}.
\end{cousinweb}

\begin{grammarlens}[title={What you've built}]
One more word for telling, explaining and describing at B1.
\end{grammarlens}

\subsection*{Your turn}
\begin{itemize}
\raggedright
  \item \emph{Say it:} \emph{merum, merī}
  \item \emph{Say it:} \emph{merum, merī}, once more
  \item \emph{Say it:} a sentence of your own with \emph{merum, merī}
  \item \emph{Recall:} say the Latin for hoarse, then the Latin for flat wine, then say \emph{merum, merī} again
\end{itemize}

\begin{tcolorbox}[breakable,colback=teal!4,colframe=teal!35!black,title={Before you move on}]
What does merum, merī mean? (\textquotedblleft{}Neat wine, wine without water\textquotedblright{}.) Say it once more.
\end{tcolorbox}
\section[dīluō, dīluere]{dīluō, dīluere — to water down, to dilute; to wash away}

\label{lesson:LA-C449-diluo}



\begin{quote}
Before the new one: say the Latin for flat wine, then the Latin for neat wine.
\end{quote}

\subsection*{You'll want to know: dīluō, dīluere}
\textbf{dīluō, dīluere} — \textquotedblleft{}to water down, to dilute; to wash away\textquotedblright{}.

\begin{quote}
Caupō vīnum aquā dīluit et magnō pretiō vendit.
\end{quote}

\begin{quote}
\emph{The innkeeper waters the wine and sells it at a high price.}
\end{quote}

\emph{dīluō, dīluere, dīluī, dīlūtum}. English \emph{dilute} comes from its participle and \emph{deluge} from the noun \emph{dīluvium}.

\begin{cousinweb}
\emph{dīluō} is built from \emph{dis-} 'apart' and \emph{luō} 'to wash', a verb akin to \emph{lavō} (Lewis \& Short). English \emph{dilute} comes from its participle \emph{dīlūtus}, and \emph{deluge} from the noun \emph{dīluvium} 'flood'.
\end{cousinweb}

\begin{grammarlens}[title={What you've built}]
One more word for telling, explaining and describing at B1.
\end{grammarlens}

\subsection*{Your turn}
\begin{itemize}
\raggedright
  \item \emph{Say it:} \emph{dīluō, dīluere}
  \item \emph{Say it:} \emph{dīluō, dīluere}, once more
  \item \emph{Say it:} a sentence of your own with \emph{dīluō, dīluere}
  \item \emph{Recall:} say the Latin for flat wine, then the Latin for neat wine, then say \emph{dīluō, dīluere} again
\end{itemize}

\begin{tcolorbox}[breakable,colback=teal!4,colframe=teal!35!black,title={Before you move on}]
What does dīluō, dīluere mean? (\textquotedblleft{}To water down, to dilute; to wash away\textquotedblright{}.) Say it once more.
\end{tcolorbox}
\section[pūtidus, pūtida, pūtidum]{pūtidus, pūtida, pūtidum — rotten, stinking}

\label{lesson:LA-C449-putidus}



\begin{quote}
Before the new one: say the Latin for neat wine, then the Latin for to water down.
\end{quote}

\subsection*{You'll want to know: pūtidus, pūtida, pūtidum}
\textbf{pūtidus, pūtida, pūtidum} — \textquotedblleft{}rotten, stinking\textquotedblright{}.

\begin{quote}
Carō in caupōnā pūtida est.
\end{quote}

\begin{quote}
\emph{The meat in the inn is rotten.}
\end{quote}

An adjective of three endings, from \emph{pūteō} 'to stink'. Also of style: affected, tiresome.

\begin{cousinweb}
\emph{pūtidus} is built on \emph{pūteō} 'to stink, to be rotten' (Lewis \& Short) with the adjective suffix \emph{-idus}, as in \emph{frīgidus}.
\end{cousinweb}

\begin{grammarlens}[title={What you've built}]
One more word for telling, explaining and describing at B1.
\end{grammarlens}

\subsection*{Your turn}
\begin{itemize}
\raggedright
  \item \emph{Say it:} \emph{pūtidus, pūtida, pūtidum}
  \item \emph{Say it:} \emph{pūtidus, pūtida, pūtidum}, once more
  \item \emph{Say it:} a sentence of your own with \emph{pūtidus, pūtida, pūtidum}
  \item \emph{Recall:} say the Latin for neat wine, then the Latin for to water down, then say \emph{pūtidus, pūtida, pūtidum} again
\end{itemize}

\begin{tcolorbox}[breakable,colback=teal!4,colframe=teal!35!black,title={Before you move on}]
What does pūtidus, pūtida, pūtidum mean? (\textquotedblleft{}Rotten, stinking\textquotedblright{}.) Say it once more.
\end{tcolorbox}
\section[cōpa, cōpae]{cōpa, cōpae — a tavern hostess, a barmaid}

\label{lesson:LA-C449-copa}



\begin{quote}
Before the new one: say the Latin for to water down, then the Latin for rotten.
\end{quote}

\subsection*{You'll want to know: cōpa, cōpae}
\textbf{cōpa, cōpae} — \textquotedblleft{}a tavern hostess, a barmaid\textquotedblright{}.

\begin{quote}
Cōpa vīnum hospitibus dat et cum eīs rīdet.
\end{quote}

\begin{quote}
\emph{The tavern hostess gives the guests wine and laughs with them.}
\end{quote}

First declension, feminine: \emph{cōpa, cōpae f.} The woman who kept or served in a tavern.

\begin{cousinweb}
\emph{cōpa} is the feminine of \emph{cōpō}, another form of \emph{caupō} 'innkeeper' (Lewis \& Short), with the \emph{au} turned to \emph{ō} as in everyday speech.
\end{cousinweb}

\begin{grammarlens}[title={What you've built}]
One more word for telling, explaining and describing at B1.
\end{grammarlens}

\subsection*{Your turn}
\begin{itemize}
\raggedright
  \item \emph{Say it:} \emph{cōpa, cōpae}
  \item \emph{Say it:} \emph{cōpa, cōpae}, once more
  \item \emph{Say it:} a sentence of your own with \emph{cōpa, cōpae}
  \item \emph{Recall:} say the Latin for to water down, then the Latin for rotten, then say \emph{cōpa, cōpae} again
\end{itemize}

\begin{tcolorbox}[breakable,colback=teal!4,colframe=teal!35!black,title={Before you move on}]
What does cōpa, cōpae mean? (\textquotedblleft{}A tavern hostess, a barmaid\textquotedblright{}.) Say it once more.
\end{tcolorbox}
